\section{The fields of the frame}
\label{sec:frame}

The course's protocol uses the following structure. The number is the length of the field in
bytes:

\begin{console}
SOF (2) | LEN (1) | TYPE (1) | DST (1) | SRC (1) | SEQ (2) | DATA (N) | CHK (2)
\end{console}

\begin{booktable}{@{}llL@{}}
\thead{Field} & \thead{Bytes} & \thead{What it solves} \\ \midrule
SOF & 2 & \emph{Where does the message begin?} A fixed pattern, \code{0xA5F7}, marking the
  start. \\
LEN & 1 & \emph{How long is it?} The number of data bytes that follow. \\
TYPE & 1 & \emph{What does it mean?} The kind of message. \\
DST & 1 & \emph{To whom?} The receiver's address. \\
SRC & 1 & \emph{From whom?} The sender's address. \\
SEQ & 2 & A sequence number, which makes it possible to pair a reply with its request and to
  recognise a duplicate. \\
DATA & N & The payload, \code{LEN} bytes long. \\
CHK & 2 & \emph{Did it arrive intact?} The sum of all preceding bytes. \\
\end{booktable}

The frame types in the course's protocol:

\begin{narrowtable}{@{}ll@{}}
\thead{Type} & \thead{Value} \\ \midrule
\code{Ping} & \code{0x00} \\
\code{Pong} & \code{0x01} \\
\code{StatusRequest} & \code{0x02} \\
\code{StatusResponse} & \code{0x03} \\
\code{Ack} & \code{0x04} \\
\code{Nack} & \code{0x05} \\
\end{narrowtable}

\subsection{Why SOF is two bytes}

One byte is not enough. Whichever value you pick will sooner or later turn up in the middle of a
payload, and then the receiver finds a frame that does not exist.

Two bytes make it unlikely rather than impossible --- \code{0xA5F7} can still occur in data. That
is why SOF on its own is never sufficient: the receiver also has to check that the length is
plausible and that the checksum is correct before it believes the frame. \Kapref{ch:parsning}
builds that check.

\subsection{Why LEN is needed when DATA runs out anyway}

Because the receiver reads one byte at a time and cannot look ahead. Without a length field it
never knows whether the next byte is data or checksum.

The alternative would be an end marker, but then that pattern has to be impossible in the payload
--- which requires re-encoding the data. CAN solves the same problem with bit stuffing (see the
course's CAN book); here a length field is enough.

\subsection{The checksum}

The course's checksum is simple: a \textbf{16-bit sum of every byte up to the CHK field}.

\begin{cppcode}
std::uint16_t computeChecksum(const std::uint8_t* buf, const std::size_t chkOffset) noexcept
{
    std::uint16_t checksum{};
    for (std::size_t i{}; i < chkOffset; ++i)
    {
        checksum += buf[i];
    }
    return checksum;
}
\end{cppcode}

It detects single bit errors and most short error bursts. What it does not detect is two bytes
swapping places, or one byte increasing by one while another decreases by one --- the sum is the
same. A CRC does better, and that is what CAN uses.

\begin{aside}
\note The checksum says that the frame arrived \emph{intact}. It says nothing about it arriving
\emph{at all}. The difference between integrity and delivery guarantee is the whole second half of
\kapref{ch:bussar}, and it is worth keeping the two apart from the start.
\end{aside}

\subsection{Serialisation and endianness}

To serialise is to write the fields of the frame into a buffer in the right order. Fields wider
than a byte --- SOF, SEQ and CHK --- have to be written in a definite byte order, and the course's
protocol uses \textbf{big-endian}: most significant byte first.

\begin{console}
SEQ = 0x7F05  ->  7F 05      (big-endian, what the protocol says)
              ->  05 7F      (little-endian, what your x86 does in memory)
\end{console}

This is the classic bug in all protocol code, and it is unpleasant for two reasons. It is
invisible when both sides have the same bug, and it looks like random corruption when they do
not: \code{0x0005} becomes \code{0x0500}, that is, 1280 instead of 5.

So never write out a multi-byte variable by copying memory. Write it byte by byte, with explicit
shifts:

\begin{cppcode}
// Big-endian: most significant byte first, whatever the host's byte order is.
buf[offset + 0U] = static_cast<std::uint8_t>(value >> 8U);
buf[offset + 1U] = static_cast<std::uint8_t>(value);
\end{cppcode}

\subsection{A complete example}

A \code{StatusRequest} from address \code{0x17} to address \code{0x25}, with sequence number
\code{0x7F05} and no payload:

\begin{console}
A5 F7 00 02 25 17 7F 05 02 5E
|---| |  |  |  |  |---| |---|
SOF  LEN TYPE      SEQ   CHK
        DST--^  ^--SRC
\end{console}

The checksum: $\mathtt{A5} + \mathtt{F7} + \mathtt{00} + \mathtt{02} + \mathtt{25} + \mathtt{17} +
\mathtt{7F} + \mathtt{05} = 606 = \mathtt{0x025E}$.

Note that \code{DST} comes before \code{SRC}. The order is arbitrary but has to be the same on
both sides, and mixing them up gives you a system where every node stubbornly answers itself.
